\scanpage{109}
la dérivation $\Theta_\delta$.

Inversement, partant de $a_1,a_2,a_3\in k$, \deleted{satisfaisant} \inserted{tels que} le polynôme
\[
 Z^3+a_1Z^2+a_2Z+a_3\in k[Z]
\]
soit sans racines multiples, i.e tels que
\[
 \underbrace{\Delta(a_1,a_2,a_3)}_{\text{discriminant}}\ne0\,, \tag*{(34)}
\]
on vérifie que les formules précédentes définissent une courbe elliptique sur $k$ (qu'on peut \deleted{définir} \inserted{construire} comme le $\operatorname{Proj}(F_*)$, où $F_*$ est défini par (31)), \deleted{\unsure{dont voici}} \inserted{en tant} que rev. \deleted{\unsure{non ramifié}} \inserted{quadratique} de $\mathbb P^1_k$ défini par
\diagram{\begin{tikzcd}[ampersand replacement=\&,column sep=small,row sep=small]
 (35)\,\deleted{\unsure{p}}\qquad p:\ \& E \arrow[rr] \arrow[dr] \& \& \mathbb P^1_k\,\deleted{\unsure{\simeq}}\\
 \& \& E/\{\pm\operatorname{id}_E\}\arrow[ur,"\unsure{\sim}"] \&
\end{tikzcd}}
ramifié en \inserted{\unsure{l'ens.|l'im.} $J$} des trois points \inserted{$z_i$} de $\mathbb P^1_{(k)}\smallsetminus\{\infty\}$ \inserted{$=k$} solutions de l'équation
\[
 Z^3+a_1Z^2+a_2Z+a_3=0\,, \tag*{(36)}
\]
\deleted{la section} \inserted{l'élément} nul\deleted{\unsure{le|e}} de $E$ étant au dessus de \deleted{la section} \inserted{l'élément} $\infty$ de $\mathbb P^1_k$. \deleted{qui est} Bien \deleted{\unsure{que}} entendu, l'inclusion $J\subset\mathbb P^1_k$ définit un isom canonique
\[
 \Sigma_J\simeq\mathbb P^1_k,\qquad t\mapsto\infty\,. \tag*{(37}
\]
% archive sheet number: 108
